% ======================================================================
% ba-grafik.tex – einheitlicher Stil für Abbildungen, Diagramme und Tabellen
% Eingebunden in main.tex per % ======================================================================
% ba-grafik.tex – einheitlicher Stil für Abbildungen, Diagramme und Tabellen
% Eingebunden in main.tex per % ======================================================================
% ba-grafik.tex – einheitlicher Stil für Abbildungen, Diagramme und Tabellen
% Eingebunden in main.tex per % ======================================================================
% ba-grafik.tex – einheitlicher Stil für Abbildungen, Diagramme und Tabellen
% Eingebunden in main.tex per \input{ba-grafik}. Regeln: CLAUDE.md, Abschnitt
% "Grafiken, Diagramme, Tabellen". Python-Gegenstück: scripts/ba_plotstil.py
% (gleiche Farben). Farben/Stile NUR hier ändern, nie in einzelnen Grafiken.
% ======================================================================
\RequirePackage{xcolor}
\RequirePackage{tikz}
\usetikzlibrary{arrows.meta, calc, positioning, angles, quotes, fit, backgrounds, shapes.geometric, babel}  % babel: "-Labels trotz ngerman
\RequirePackage{pgfplots}
\pgfplotsset{compat=1.18}
\RequirePackage{tabularx}   % Tabellen auf Textbreite (Spaltentyp X)
\RequirePackage{xltabular}  % mehrseitige Tabellen mit wiederholtem Kopf (\endhead)
% booktabs lädt main.tex bereits

% ---------------------------------------------------------------------
% 1  Farben  (FAPS-Grün aus main.tex als Hauptfarbe)
%    Reihenfolge = Reihenfolge der Datenreihen in Diagrammen.
%    baAkzent nur für das, was hervorgehoben werden soll (gesucht,
%    unbekannt, Fehler, Grenzwert) – nie als normale Datenreihe.
% ---------------------------------------------------------------------
\colorlet{baGruen}{FAPSGruen}                 % RGB 151,193,57  – Hauptfarbe, Reihe 1
\definecolor{baBlau}{RGB}{4, 49, 106}         % #04316A – FAU-Blau, Reihe 2
\definecolor{baGrau}{RGB}{89, 89, 89}         % #595959 – Reihe 3, Text in Grafiken
\definecolor{baHellblau}{RGB}{114, 163, 205}  % #72A3CD – Reihe 4
\definecolor{baAkzent}{RGB}{214, 96, 0}       % #D66000 – Hervorhebung (gesucht/kritisch)
\definecolor{baHellgrau}{RGB}{217, 217, 217}  % #D9D9D9 – Hilfslinien, Gitter, Flächen
\colorlet{baFlaeche}{FAPSGruen!15}            % gefüllte Blöcke (nur wenn Farbe Bedeutung trägt)

% ---------------------------------------------------------------------
% 2  TikZ-Grundstil  (Schrift = Fließtextschrift, \small; Linienstärken fest)
% ---------------------------------------------------------------------
\tikzset{
  ba/.style={font=\small, line cap=round, line join=round, >={Stealth[length=2.2mm]}},
  % Linien
  ba duenn/.style={line width=0.4pt},
  ba normal/.style={line width=0.8pt},
  ba dick/.style={line width=1.2pt},
  ba hilfslinie/.style={ba normal, dash pattern=on 3pt off 2pt, color=baGrau},
  ba pfeil/.style={ba normal, ->},
  ba doppelpfeil/.style={ba normal, <->},
  % Blockdiagramme / Abläufe
  ba block/.style={draw=baGrau, ba normal, fill=white, rounded corners=1.5pt,
                   minimum height=8mm, minimum width=22mm, align=center, inner sep=3pt},
  ba block hervor/.style={ba block, draw=baGruen, fill=baFlaeche},
  ba gruppe/.style={draw=baHellgrau, ba normal, dash pattern=on 3pt off 2pt, rounded corners=3pt, inner sep=5pt},
  % Kinematik: Glieder, Gelenke, Endeffektor, Winkel
  ba glied/.style={line width=2.4pt, color=baGruen},
  ba gelenk/.style={circle, fill=baGruen, draw=white, line width=0.6pt, minimum size=3.2mm, inner sep=0pt},
  ba basis/.style={ba gelenk, fill=baGrau},
  ba endeffektor/.style={circle, draw=baGruen, fill=white, line width=1.2pt, minimum size=3.6mm, inner sep=0pt},
  ba winkel/.style={draw, ba normal, ->, angle radius=9mm, angle eccentricity=1.4},
  % Hervorhebung von Gesuchtem/Unbekanntem (Farbe + Text, nie Farbe allein)
  ba gesucht/.style={color=baAkzent, text=baAkzent},
  % Bild in mehreren Teilen nebeneinander: Trennlinie und Teilbeschriftung
  ba trenner/.style={draw=baHellgrau, ba normal},
  ba teiltitel/.style={font=\small, text=black, anchor=north},
}

% ---------------------------------------------------------------------
% 3  pgfplots-Grundstil  (Achsen "Größe in Einheit", Dezimalkomma, Gitter)
%    Verwendung: \begin{axis}[ba plot, xlabel={Zeit in s}, ylabel={...}]
% ---------------------------------------------------------------------
\pgfplotscreateplotcyclelist{ba farben}{%
  {baGruen, mark=*, mark size=1.6pt, line width=1pt},
  {baBlau, mark=square*, mark size=1.5pt, line width=1pt},
  {baGrau, mark=triangle*, mark size=1.9pt, line width=1pt},
  {baHellblau, mark=diamond*, mark size=1.9pt, line width=1pt}%
}
\pgfplotsset{
  ba plot/.style={
    width=\linewidth, height=0.5\linewidth,
    scale only axis=false,
    axis lines=box,                               % Rahmen um das Diagramm
    axis line style={baGrau, line width=0.6pt},
    tick style={baGrau, line width=0.6pt},
    tick align=inside,
    every tick label/.append style={font=\small},
    /pgf/number format/assume math mode=true,     % Zahlen in Textschrift (Arial), nicht in Mathe-Serifen
    label style={font=\small},
    title style={font=\small},
    grid=major,
    major grid style={baHellgrau, line width=0.4pt},
    cycle list name=ba farben,
    legend style={font=\small, draw=baHellgrau, fill=white, fill opacity=0.9,
                  text opacity=1, rounded corners=1pt, cells={anchor=west}},
    legend cell align=left,
    /pgf/number format/.cd, use comma, 1000 sep={.},   % 1.350 und 1,5
  },
  % Linien ohne Marker (z. B. dichte Messreihen)
  ba ohne marker/.style={no markers},
  % Balkendiagramm
  ba balken/.style={ba plot, ybar, bar width=10pt, enlarge x limits=0.15,
                    every axis plot/.append style={fill, draw=none}},
}
% Grenzwert/Sollwert im Diagramm: \draw[ba grenzwert] (axis cs:0,1.5) -- (axis cs:4,1.5);
\tikzset{ba grenzwert/.style={baAkzent, line width=0.8pt, dash pattern=on 4pt off 2pt}}

% ---------------------------------------------------------------------
% 4  Tabellen – zwei Stile, je nach Inhalt wählen (Regel in CLAUDE.md):
%    a) Gitter  (Standard für Text-/Vergleichstabellen): senkrechte und
%       waagrechte Linien, Kopfzeile hellgrün:  |L|L|  + \hline + \tabkopfzeile
%    b) Schlicht (für reine Zahlen-/Messwerttabellen): booktabs
%       \toprule/\midrule/\bottomrule, keine senkrechten Linien
% ---------------------------------------------------------------------
\RequirePackage{colortbl}    % Linienfarbe und hinterlegte Kopfzeile
\RequirePackage{array}
% Linksbündige X-Spalte ohne Blocksatz (Richtlinien: kein Blocksatz in Tabellen)
% (\hspace{0pt} erlaubt Silbentrennung auch im ersten Wort einer Zelle)
\newcolumntype{L}{>{\raggedright\arraybackslash\hspace{0pt}}X}
\newcolumntype{C}{>{\centering\arraybackslash\hspace{0pt}}X}
% Kopfzelle fett: \tabkopf{Spalte};  Kopfzeile hinterlegt: \tabkopfzeile vor die erste Zeile
\newcommand{\tabkopf}[1]{\textbf{#1}}
\newcommand{\tabkopfzeile}{\rowcolor{baFlaeche}}
% Opt-in je Tabelle: \tabstil als erste Zeile in der table-Umgebung (vor \caption ist egal,
% die Caption formatiert weiterhin das caption-Paket der Vorlage). Ändert nichts global.
\newcommand{\tabstil}{\small\renewcommand{\arraystretch}{1.15}%
  \setlength{\extrarowheight}{1.5pt}\setlength{\arrayrulewidth}{0.5pt}\arrayrulecolor{baGrau}}

% ---------------------------------------------------------------------
% Hinweis: Diese Datei ändert NICHTS an der FAPS-Vorlage (Captions, Überschriften,
% Verzeichnisse, Nummerierung, Seitenlayout). Sie definiert nur Farben, Stile und
% Befehle, die erst wirken, wenn eine Grafik/Tabelle sie ausdrücklich benutzt.
% ---------------------------------------------------------------------
. Regeln: CLAUDE.md, Abschnitt
% "Grafiken, Diagramme, Tabellen". Python-Gegenstück: scripts/ba_plotstil.py
% (gleiche Farben). Farben/Stile NUR hier ändern, nie in einzelnen Grafiken.
% ======================================================================
\RequirePackage{xcolor}
\RequirePackage{tikz}
\usetikzlibrary{arrows.meta, calc, positioning, angles, quotes, fit, backgrounds, shapes.geometric, babel}  % babel: "-Labels trotz ngerman
\RequirePackage{pgfplots}
\pgfplotsset{compat=1.18}
\RequirePackage{tabularx}   % Tabellen auf Textbreite (Spaltentyp X)
\RequirePackage{xltabular}  % mehrseitige Tabellen mit wiederholtem Kopf (\endhead)
% booktabs lädt main.tex bereits

% ---------------------------------------------------------------------
% 1  Farben  (FAPS-Grün aus main.tex als Hauptfarbe)
%    Reihenfolge = Reihenfolge der Datenreihen in Diagrammen.
%    baAkzent nur für das, was hervorgehoben werden soll (gesucht,
%    unbekannt, Fehler, Grenzwert) – nie als normale Datenreihe.
% ---------------------------------------------------------------------
\colorlet{baGruen}{FAPSGruen}                 % RGB 151,193,57  – Hauptfarbe, Reihe 1
\definecolor{baBlau}{RGB}{4, 49, 106}         % #04316A – FAU-Blau, Reihe 2
\definecolor{baGrau}{RGB}{89, 89, 89}         % #595959 – Reihe 3, Text in Grafiken
\definecolor{baHellblau}{RGB}{114, 163, 205}  % #72A3CD – Reihe 4
\definecolor{baAkzent}{RGB}{214, 96, 0}       % #D66000 – Hervorhebung (gesucht/kritisch)
\definecolor{baHellgrau}{RGB}{217, 217, 217}  % #D9D9D9 – Hilfslinien, Gitter, Flächen
\colorlet{baFlaeche}{FAPSGruen!15}            % gefüllte Blöcke (nur wenn Farbe Bedeutung trägt)

% ---------------------------------------------------------------------
% 2  TikZ-Grundstil  (Schrift = Fließtextschrift, \small; Linienstärken fest)
% ---------------------------------------------------------------------
\tikzset{
  ba/.style={font=\small, line cap=round, line join=round, >={Stealth[length=2.2mm]}},
  % Linien
  ba duenn/.style={line width=0.4pt},
  ba normal/.style={line width=0.8pt},
  ba dick/.style={line width=1.2pt},
  ba hilfslinie/.style={ba normal, dash pattern=on 3pt off 2pt, color=baGrau},
  ba pfeil/.style={ba normal, ->},
  ba doppelpfeil/.style={ba normal, <->},
  % Blockdiagramme / Abläufe
  ba block/.style={draw=baGrau, ba normal, fill=white, rounded corners=1.5pt,
                   minimum height=8mm, minimum width=22mm, align=center, inner sep=3pt},
  ba block hervor/.style={ba block, draw=baGruen, fill=baFlaeche},
  ba gruppe/.style={draw=baHellgrau, ba normal, dash pattern=on 3pt off 2pt, rounded corners=3pt, inner sep=5pt},
  % Kinematik: Glieder, Gelenke, Endeffektor, Winkel
  ba glied/.style={line width=2.4pt, color=baGruen},
  ba gelenk/.style={circle, fill=baGruen, draw=white, line width=0.6pt, minimum size=3.2mm, inner sep=0pt},
  ba basis/.style={ba gelenk, fill=baGrau},
  ba endeffektor/.style={circle, draw=baGruen, fill=white, line width=1.2pt, minimum size=3.6mm, inner sep=0pt},
  ba winkel/.style={draw, ba normal, ->, angle radius=9mm, angle eccentricity=1.4},
  % Hervorhebung von Gesuchtem/Unbekanntem (Farbe + Text, nie Farbe allein)
  ba gesucht/.style={color=baAkzent, text=baAkzent},
  % Bild in mehreren Teilen nebeneinander: Trennlinie und Teilbeschriftung
  ba trenner/.style={draw=baHellgrau, ba normal},
  ba teiltitel/.style={font=\small, text=black, anchor=north},
}

% ---------------------------------------------------------------------
% 3  pgfplots-Grundstil  (Achsen "Größe in Einheit", Dezimalkomma, Gitter)
%    Verwendung: \begin{axis}[ba plot, xlabel={Zeit in s}, ylabel={...}]
% ---------------------------------------------------------------------
\pgfplotscreateplotcyclelist{ba farben}{%
  {baGruen, mark=*, mark size=1.6pt, line width=1pt},
  {baBlau, mark=square*, mark size=1.5pt, line width=1pt},
  {baGrau, mark=triangle*, mark size=1.9pt, line width=1pt},
  {baHellblau, mark=diamond*, mark size=1.9pt, line width=1pt}%
}
\pgfplotsset{
  ba plot/.style={
    width=\linewidth, height=0.5\linewidth,
    scale only axis=false,
    axis lines=box,                               % Rahmen um das Diagramm
    axis line style={baGrau, line width=0.6pt},
    tick style={baGrau, line width=0.6pt},
    tick align=inside,
    every tick label/.append style={font=\small},
    /pgf/number format/assume math mode=true,     % Zahlen in Textschrift (Arial), nicht in Mathe-Serifen
    label style={font=\small},
    title style={font=\small},
    grid=major,
    major grid style={baHellgrau, line width=0.4pt},
    cycle list name=ba farben,
    legend style={font=\small, draw=baHellgrau, fill=white, fill opacity=0.9,
                  text opacity=1, rounded corners=1pt, cells={anchor=west}},
    legend cell align=left,
    /pgf/number format/.cd, use comma, 1000 sep={.},   % 1.350 und 1,5
  },
  % Linien ohne Marker (z. B. dichte Messreihen)
  ba ohne marker/.style={no markers},
  % Balkendiagramm
  ba balken/.style={ba plot, ybar, bar width=10pt, enlarge x limits=0.15,
                    every axis plot/.append style={fill, draw=none}},
}
% Grenzwert/Sollwert im Diagramm: \draw[ba grenzwert] (axis cs:0,1.5) -- (axis cs:4,1.5);
\tikzset{ba grenzwert/.style={baAkzent, line width=0.8pt, dash pattern=on 4pt off 2pt}}

% ---------------------------------------------------------------------
% 4  Tabellen – zwei Stile, je nach Inhalt wählen (Regel in CLAUDE.md):
%    a) Gitter  (Standard für Text-/Vergleichstabellen): senkrechte und
%       waagrechte Linien, Kopfzeile hellgrün:  |L|L|  + \hline + \tabkopfzeile
%    b) Schlicht (für reine Zahlen-/Messwerttabellen): booktabs
%       \toprule/\midrule/\bottomrule, keine senkrechten Linien
% ---------------------------------------------------------------------
\RequirePackage{colortbl}    % Linienfarbe und hinterlegte Kopfzeile
\RequirePackage{array}
% Linksbündige X-Spalte ohne Blocksatz (Richtlinien: kein Blocksatz in Tabellen)
% (\hspace{0pt} erlaubt Silbentrennung auch im ersten Wort einer Zelle)
\newcolumntype{L}{>{\raggedright\arraybackslash\hspace{0pt}}X}
\newcolumntype{C}{>{\centering\arraybackslash\hspace{0pt}}X}
% Kopfzelle fett: \tabkopf{Spalte};  Kopfzeile hinterlegt: \tabkopfzeile vor die erste Zeile
\newcommand{\tabkopf}[1]{\textbf{#1}}
\newcommand{\tabkopfzeile}{\rowcolor{baFlaeche}}
% Opt-in je Tabelle: \tabstil als erste Zeile in der table-Umgebung (vor \caption ist egal,
% die Caption formatiert weiterhin das caption-Paket der Vorlage). Ändert nichts global.
\newcommand{\tabstil}{\small\renewcommand{\arraystretch}{1.15}%
  \setlength{\extrarowheight}{1.5pt}\setlength{\arrayrulewidth}{0.5pt}\arrayrulecolor{baGrau}}

% ---------------------------------------------------------------------
% Hinweis: Diese Datei ändert NICHTS an der FAPS-Vorlage (Captions, Überschriften,
% Verzeichnisse, Nummerierung, Seitenlayout). Sie definiert nur Farben, Stile und
% Befehle, die erst wirken, wenn eine Grafik/Tabelle sie ausdrücklich benutzt.
% ---------------------------------------------------------------------
. Regeln: CLAUDE.md, Abschnitt
% "Grafiken, Diagramme, Tabellen". Python-Gegenstück: scripts/ba_plotstil.py
% (gleiche Farben). Farben/Stile NUR hier ändern, nie in einzelnen Grafiken.
% ======================================================================
\RequirePackage{xcolor}
\RequirePackage{tikz}
\usetikzlibrary{arrows.meta, calc, positioning, angles, quotes, fit, backgrounds, shapes.geometric, babel}  % babel: "-Labels trotz ngerman
\RequirePackage{pgfplots}
\pgfplotsset{compat=1.18}
\RequirePackage{tabularx}   % Tabellen auf Textbreite (Spaltentyp X)
\RequirePackage{xltabular}  % mehrseitige Tabellen mit wiederholtem Kopf (\endhead)
% booktabs lädt main.tex bereits

% ---------------------------------------------------------------------
% 1  Farben  (FAPS-Grün aus main.tex als Hauptfarbe)
%    Reihenfolge = Reihenfolge der Datenreihen in Diagrammen.
%    baAkzent nur für das, was hervorgehoben werden soll (gesucht,
%    unbekannt, Fehler, Grenzwert) – nie als normale Datenreihe.
% ---------------------------------------------------------------------
\colorlet{baGruen}{FAPSGruen}                 % RGB 151,193,57  – Hauptfarbe, Reihe 1
\definecolor{baBlau}{RGB}{4, 49, 106}         % #04316A – FAU-Blau, Reihe 2
\definecolor{baGrau}{RGB}{89, 89, 89}         % #595959 – Reihe 3, Text in Grafiken
\definecolor{baHellblau}{RGB}{114, 163, 205}  % #72A3CD – Reihe 4
\definecolor{baAkzent}{RGB}{214, 96, 0}       % #D66000 – Hervorhebung (gesucht/kritisch)
\definecolor{baHellgrau}{RGB}{217, 217, 217}  % #D9D9D9 – Hilfslinien, Gitter, Flächen
\colorlet{baFlaeche}{FAPSGruen!15}            % gefüllte Blöcke (nur wenn Farbe Bedeutung trägt)

% ---------------------------------------------------------------------
% 2  TikZ-Grundstil  (Schrift = Fließtextschrift, \small; Linienstärken fest)
% ---------------------------------------------------------------------
\tikzset{
  ba/.style={font=\small, line cap=round, line join=round, >={Stealth[length=2.2mm]}},
  % Linien
  ba duenn/.style={line width=0.4pt},
  ba normal/.style={line width=0.8pt},
  ba dick/.style={line width=1.2pt},
  ba hilfslinie/.style={ba normal, dash pattern=on 3pt off 2pt, color=baGrau},
  ba pfeil/.style={ba normal, ->},
  ba doppelpfeil/.style={ba normal, <->},
  % Blockdiagramme / Abläufe
  ba block/.style={draw=baGrau, ba normal, fill=white, rounded corners=1.5pt,
                   minimum height=8mm, minimum width=22mm, align=center, inner sep=3pt},
  ba block hervor/.style={ba block, draw=baGruen, fill=baFlaeche},
  ba gruppe/.style={draw=baHellgrau, ba normal, dash pattern=on 3pt off 2pt, rounded corners=3pt, inner sep=5pt},
  % Kinematik: Glieder, Gelenke, Endeffektor, Winkel
  ba glied/.style={line width=2.4pt, color=baGruen},
  ba gelenk/.style={circle, fill=baGruen, draw=white, line width=0.6pt, minimum size=3.2mm, inner sep=0pt},
  ba basis/.style={ba gelenk, fill=baGrau},
  ba endeffektor/.style={circle, draw=baGruen, fill=white, line width=1.2pt, minimum size=3.6mm, inner sep=0pt},
  ba winkel/.style={draw, ba normal, ->, angle radius=9mm, angle eccentricity=1.4},
  % Hervorhebung von Gesuchtem/Unbekanntem (Farbe + Text, nie Farbe allein)
  ba gesucht/.style={color=baAkzent, text=baAkzent},
  % Bild in mehreren Teilen nebeneinander: Trennlinie und Teilbeschriftung
  ba trenner/.style={draw=baHellgrau, ba normal},
  ba teiltitel/.style={font=\small, text=black, anchor=north},
}

% ---------------------------------------------------------------------
% 3  pgfplots-Grundstil  (Achsen "Größe in Einheit", Dezimalkomma, Gitter)
%    Verwendung: \begin{axis}[ba plot, xlabel={Zeit in s}, ylabel={...}]
% ---------------------------------------------------------------------
\pgfplotscreateplotcyclelist{ba farben}{%
  {baGruen, mark=*, mark size=1.6pt, line width=1pt},
  {baBlau, mark=square*, mark size=1.5pt, line width=1pt},
  {baGrau, mark=triangle*, mark size=1.9pt, line width=1pt},
  {baHellblau, mark=diamond*, mark size=1.9pt, line width=1pt}%
}
\pgfplotsset{
  ba plot/.style={
    width=\linewidth, height=0.5\linewidth,
    scale only axis=false,
    axis lines=box,                               % Rahmen um das Diagramm
    axis line style={baGrau, line width=0.6pt},
    tick style={baGrau, line width=0.6pt},
    tick align=inside,
    every tick label/.append style={font=\small},
    /pgf/number format/assume math mode=true,     % Zahlen in Textschrift (Arial), nicht in Mathe-Serifen
    label style={font=\small},
    title style={font=\small},
    grid=major,
    major grid style={baHellgrau, line width=0.4pt},
    cycle list name=ba farben,
    legend style={font=\small, draw=baHellgrau, fill=white, fill opacity=0.9,
                  text opacity=1, rounded corners=1pt, cells={anchor=west}},
    legend cell align=left,
    /pgf/number format/.cd, use comma, 1000 sep={.},   % 1.350 und 1,5
  },
  % Linien ohne Marker (z. B. dichte Messreihen)
  ba ohne marker/.style={no markers},
  % Balkendiagramm
  ba balken/.style={ba plot, ybar, bar width=10pt, enlarge x limits=0.15,
                    every axis plot/.append style={fill, draw=none}},
}
% Grenzwert/Sollwert im Diagramm: \draw[ba grenzwert] (axis cs:0,1.5) -- (axis cs:4,1.5);
\tikzset{ba grenzwert/.style={baAkzent, line width=0.8pt, dash pattern=on 4pt off 2pt}}

% ---------------------------------------------------------------------
% 4  Tabellen – zwei Stile, je nach Inhalt wählen (Regel in CLAUDE.md):
%    a) Gitter  (Standard für Text-/Vergleichstabellen): senkrechte und
%       waagrechte Linien, Kopfzeile hellgrün:  |L|L|  + \hline + \tabkopfzeile
%    b) Schlicht (für reine Zahlen-/Messwerttabellen): booktabs
%       \toprule/\midrule/\bottomrule, keine senkrechten Linien
% ---------------------------------------------------------------------
\RequirePackage{colortbl}    % Linienfarbe und hinterlegte Kopfzeile
\RequirePackage{array}
% Linksbündige X-Spalte ohne Blocksatz (Richtlinien: kein Blocksatz in Tabellen)
% (\hspace{0pt} erlaubt Silbentrennung auch im ersten Wort einer Zelle)
\newcolumntype{L}{>{\raggedright\arraybackslash\hspace{0pt}}X}
\newcolumntype{C}{>{\centering\arraybackslash\hspace{0pt}}X}
% Kopfzelle fett: \tabkopf{Spalte};  Kopfzeile hinterlegt: \tabkopfzeile vor die erste Zeile
\newcommand{\tabkopf}[1]{\textbf{#1}}
\newcommand{\tabkopfzeile}{\rowcolor{baFlaeche}}
% Opt-in je Tabelle: \tabstil als erste Zeile in der table-Umgebung (vor \caption ist egal,
% die Caption formatiert weiterhin das caption-Paket der Vorlage). Ändert nichts global.
\newcommand{\tabstil}{\small\renewcommand{\arraystretch}{1.15}%
  \setlength{\extrarowheight}{1.5pt}\setlength{\arrayrulewidth}{0.5pt}\arrayrulecolor{baGrau}}

% ---------------------------------------------------------------------
% Hinweis: Diese Datei ändert NICHTS an der FAPS-Vorlage (Captions, Überschriften,
% Verzeichnisse, Nummerierung, Seitenlayout). Sie definiert nur Farben, Stile und
% Befehle, die erst wirken, wenn eine Grafik/Tabelle sie ausdrücklich benutzt.
% ---------------------------------------------------------------------
. Regeln: CLAUDE.md, Abschnitt
% "Grafiken, Diagramme, Tabellen". Python-Gegenstück: scripts/ba_plotstil.py
% (gleiche Farben). Farben/Stile NUR hier ändern, nie in einzelnen Grafiken.
% ======================================================================
\RequirePackage{xcolor}
\RequirePackage{tikz}
\usetikzlibrary{arrows.meta, calc, positioning, angles, quotes, fit, backgrounds, shapes.geometric, babel}  % babel: "-Labels trotz ngerman
\RequirePackage{pgfplots}
\pgfplotsset{compat=1.18}
\RequirePackage{tabularx}   % Tabellen auf Textbreite (Spaltentyp X)
\RequirePackage{xltabular}  % mehrseitige Tabellen mit wiederholtem Kopf (\endhead)
% booktabs lädt main.tex bereits

% ---------------------------------------------------------------------
% 1  Farben  (FAPS-Grün aus main.tex als Hauptfarbe)
%    Reihenfolge = Reihenfolge der Datenreihen in Diagrammen.
%    baAkzent nur für das, was hervorgehoben werden soll (gesucht,
%    unbekannt, Fehler, Grenzwert) – nie als normale Datenreihe.
% ---------------------------------------------------------------------
\colorlet{baGruen}{FAPSGruen}                 % RGB 151,193,57  – Hauptfarbe, Reihe 1
\definecolor{baBlau}{RGB}{4, 49, 106}         % #04316A – FAU-Blau, Reihe 2
\definecolor{baGrau}{RGB}{89, 89, 89}         % #595959 – Reihe 3, Text in Grafiken
\definecolor{baHellblau}{RGB}{114, 163, 205}  % #72A3CD – Reihe 4
\definecolor{baAkzent}{RGB}{214, 96, 0}       % #D66000 – Hervorhebung (gesucht/kritisch)
\definecolor{baHellgrau}{RGB}{217, 217, 217}  % #D9D9D9 – Hilfslinien, Gitter, Flächen
\colorlet{baFlaeche}{FAPSGruen!15}            % gefüllte Blöcke (nur wenn Farbe Bedeutung trägt)

% ---------------------------------------------------------------------
% 2  TikZ-Grundstil  (Schrift = Fließtextschrift, \small; Linienstärken fest)
% ---------------------------------------------------------------------
\tikzset{
  ba/.style={font=\small, line cap=round, line join=round, >={Stealth[length=2.2mm]}},
  % Linien
  ba duenn/.style={line width=0.4pt},
  ba normal/.style={line width=0.8pt},
  ba dick/.style={line width=1.2pt},
  ba hilfslinie/.style={ba normal, dash pattern=on 3pt off 2pt, color=baGrau},
  ba pfeil/.style={ba normal, ->},
  ba doppelpfeil/.style={ba normal, <->},
  % Blockdiagramme / Abläufe
  ba block/.style={draw=baGrau, ba normal, fill=white, rounded corners=1.5pt,
                   minimum height=8mm, minimum width=22mm, align=center, inner sep=3pt},
  ba block hervor/.style={ba block, draw=baGruen, fill=baFlaeche},
  ba gruppe/.style={draw=baHellgrau, ba normal, dash pattern=on 3pt off 2pt, rounded corners=3pt, inner sep=5pt},
  % Kinematik: Glieder, Gelenke, Endeffektor, Winkel
  ba glied/.style={line width=2.4pt, color=baGruen},
  ba gelenk/.style={circle, fill=baGruen, draw=white, line width=0.6pt, minimum size=3.2mm, inner sep=0pt},
  ba basis/.style={ba gelenk, fill=baGrau},
  ba endeffektor/.style={circle, draw=baGruen, fill=white, line width=1.2pt, minimum size=3.6mm, inner sep=0pt},
  ba winkel/.style={draw, ba normal, ->, angle radius=9mm, angle eccentricity=1.4},
  % Hervorhebung von Gesuchtem/Unbekanntem (Farbe + Text, nie Farbe allein)
  ba gesucht/.style={color=baAkzent, text=baAkzent},
  % Bild in mehreren Teilen nebeneinander: Trennlinie und Teilbeschriftung
  ba trenner/.style={draw=baHellgrau, ba normal},
  ba teiltitel/.style={font=\small, text=black, anchor=north},
}

% ---------------------------------------------------------------------
% 3  pgfplots-Grundstil  (Achsen "Größe in Einheit", Dezimalkomma, Gitter)
%    Verwendung: \begin{axis}[ba plot, xlabel={Zeit in s}, ylabel={...}]
% ---------------------------------------------------------------------
\pgfplotscreateplotcyclelist{ba farben}{%
  {baGruen, mark=*, mark size=1.6pt, line width=1pt},
  {baBlau, mark=square*, mark size=1.5pt, line width=1pt},
  {baGrau, mark=triangle*, mark size=1.9pt, line width=1pt},
  {baHellblau, mark=diamond*, mark size=1.9pt, line width=1pt}%
}
\pgfplotsset{
  ba plot/.style={
    width=\linewidth, height=0.5\linewidth,
    scale only axis=false,
    axis lines=box,                               % Rahmen um das Diagramm
    axis line style={baGrau, line width=0.6pt},
    tick style={baGrau, line width=0.6pt},
    tick align=inside,
    every tick label/.append style={font=\small},
    /pgf/number format/assume math mode=true,     % Zahlen in Textschrift (Arial), nicht in Mathe-Serifen
    label style={font=\small},
    title style={font=\small},
    grid=major,
    major grid style={baHellgrau, line width=0.4pt},
    cycle list name=ba farben,
    legend style={font=\small, draw=baHellgrau, fill=white, fill opacity=0.9,
                  text opacity=1, rounded corners=1pt, cells={anchor=west}},
    legend cell align=left,
    /pgf/number format/.cd, use comma, 1000 sep={.},   % 1.350 und 1,5
  },
  % Linien ohne Marker (z. B. dichte Messreihen)
  ba ohne marker/.style={no markers},
  % Balkendiagramm
  ba balken/.style={ba plot, ybar, bar width=10pt, enlarge x limits=0.15,
                    every axis plot/.append style={fill, draw=none}},
}
% Grenzwert/Sollwert im Diagramm: \draw[ba grenzwert] (axis cs:0,1.5) -- (axis cs:4,1.5);
\tikzset{ba grenzwert/.style={baAkzent, line width=0.8pt, dash pattern=on 4pt off 2pt}}

% ---------------------------------------------------------------------
% 4  Tabellen – zwei Stile, je nach Inhalt wählen (Regel in CLAUDE.md):
%    a) Gitter  (Standard für Text-/Vergleichstabellen): senkrechte und
%       waagrechte Linien, Kopfzeile hellgrün:  |L|L|  + \hline + \tabkopfzeile
%    b) Schlicht (für reine Zahlen-/Messwerttabellen): booktabs
%       \toprule/\midrule/\bottomrule, keine senkrechten Linien
% ---------------------------------------------------------------------
\RequirePackage{colortbl}    % Linienfarbe und hinterlegte Kopfzeile
\RequirePackage{array}
% Linksbündige X-Spalte ohne Blocksatz (Richtlinien: kein Blocksatz in Tabellen)
% (\hspace{0pt} erlaubt Silbentrennung auch im ersten Wort einer Zelle)
\newcolumntype{L}{>{\raggedright\arraybackslash\hspace{0pt}}X}
\newcolumntype{C}{>{\centering\arraybackslash\hspace{0pt}}X}
% Kopfzelle fett: \tabkopf{Spalte};  Kopfzeile hinterlegt: \tabkopfzeile vor die erste Zeile
\newcommand{\tabkopf}[1]{\textbf{#1}}
\newcommand{\tabkopfzeile}{\rowcolor{baFlaeche}}
% Opt-in je Tabelle: \tabstil als erste Zeile in der table-Umgebung (vor \caption ist egal,
% die Caption formatiert weiterhin das caption-Paket der Vorlage). Ändert nichts global.
\newcommand{\tabstil}{\small\renewcommand{\arraystretch}{1.15}%
  \setlength{\extrarowheight}{1.5pt}\setlength{\arrayrulewidth}{0.5pt}\arrayrulecolor{baGrau}}

% ---------------------------------------------------------------------
% Hinweis: Diese Datei ändert NICHTS an der FAPS-Vorlage (Captions, Überschriften,
% Verzeichnisse, Nummerierung, Seitenlayout). Sie definiert nur Farben, Stile und
% Befehle, die erst wirken, wenn eine Grafik/Tabelle sie ausdrücklich benutzt.
% ---------------------------------------------------------------------
